% GENERATED FILE. Edit canonical lessons, then run npm run generate:books.
% canonical-source-hash: fnv1a64:905b831d6f0795fb
% canonical-lessons: GU-C472-nas, GU-C472-snayu, GU-C472-khopri, GU-C472-atardu, GU-C472-hojri

\chapter{Inside the Body}
\label{ch:gu-inside-the-body}
\hlchaptermodality{472}

\begin{chapteropening}
\textbf{By the end of this chapter:} \emph{I can name some parts inside the body.}

Complete the last lesson of chapter 472: I can name some parts inside the body.
\end{chapteropening}

\section[\texorpdfstring{\gu{નસ} (nas)}{નસ (nas)}]{\gu{નસ} (nas) — a vein; a nerve}

\label{lesson:GU-C472-nas}



\begin{quote}
Before the new one: say the Gujarati for wrinkles, then the Gujarati for grey hairs.
\end{quote}

\subsection*{You'll want to know: \gu{નસ}}
\textbf{\gu{નસ}} — \emph{nas} — \textquotedblleft{}a vein; a nerve\textquotedblright{}.

\begin{quote}
\gu{દાક્તરે} \gu{મારા} \gu{હાથની} \gu{નસ} \gu{તપાસી}.
\end{quote}

\begin{quote}
\emph{dāktare mārā hāthnī nas tapāsī.}
\end{quote}

\begin{quote}
The doctor examined the vein in my arm.
\end{quote}

A feminine noun.

\begin{cousinweb}
\textbf{\gu{નસ}} (\emph{nas}) is inherited from Sanskrit \emph{snasā}, \textquotedblleft{}tendon, muscle\textquotedblright{}; CDIAL lists the Gujarati word, \textquotedblleft{}vein, nerve, tendon\textquotedblright{} (CDIAL 13784).
\end{cousinweb}

\begin{grammarlens}[title={What you've built}]
One more word for describing what you lived and felt.
\end{grammarlens}

\subsection*{Your turn}
\begin{itemize}
\raggedright
  \item \emph{Say it:} \emph{nas}
  \item \emph{Say it:} \emph{nas}, once more
  \item \emph{Say it:} \emph{dāktare mārā hāthnī nas tapāsī.} — The doctor examined the vein in my arm.
  \item \emph{Recall:} say the Gujarati for wrinkles, then the Gujarati for grey hairs, then say \emph{nas} again
\end{itemize}

\begin{tcolorbox}[breakable,colback=teal!4,colframe=teal!35!black,title={Before you move on}]
What does \gu{નસ} mean? (\textquotedblleft{}A vein; a nerve\textquotedblright{}.) Say it once more.
\end{tcolorbox}
\section[\texorpdfstring{\gu{સ્નાયુ} (snāyu)}{સ્નાયુ (snāyu)}]{\gu{સ્નાયુ} (snāyu) — a muscle}

\label{lesson:GU-C472-snayu}



\begin{quote}
Before the new one: say the Gujarati for grey hairs, then the Gujarati for a vein.
\end{quote}

\subsection*{You'll want to know: \gu{સ્નાયુ}}
\textbf{\gu{સ્નાયુ}} — \emph{snāyu} — \textquotedblleft{}a muscle\textquotedblright{}.

\begin{quote}
\gu{કસરતથી} \gu{સ્નાયુ} \gu{મજબૂત} \gu{બને} \gu{છે}.
\end{quote}

\begin{quote}
\emph{kasratthī snāyu majbūt bane che.}
\end{quote}

\begin{quote}
Exercise makes the muscles strong.
\end{quote}

A masculine noun.

\begin{cousinweb}
\textbf{\gu{સ્નાયુ}} (\emph{snāyu}) is a tatsama, Sanskrit \emph{snāyu}, which Monier-Williams gives as \textquotedblleft{}any sinew or ligament\textquotedblright{}, \textquotedblleft{}tendon, muscle, nerve\textquotedblright{}; modern Gujarati uses it for a muscle.
\end{cousinweb}

\begin{grammarlens}[title={What you've built}]
One more word for describing what you lived and felt.
\end{grammarlens}

\subsection*{Your turn}
\begin{itemize}
\raggedright
  \item \emph{Say it:} \emph{snāyu}
  \item \emph{Say it:} \emph{snāyu}, once more
  \item \emph{Say it:} \emph{kasratthī snāyu majbūt bane che.} — Exercise makes the muscles strong.
  \item \emph{Recall:} say the Gujarati for grey hairs, then the Gujarati for a vein, then say \emph{snāyu} again
\end{itemize}

\begin{tcolorbox}[breakable,colback=teal!4,colframe=teal!35!black,title={Before you move on}]
What does \gu{સ્નાયુ} mean? (\textquotedblleft{}A muscle\textquotedblright{}.) Say it once more.
\end{tcolorbox}
\section[\texorpdfstring{\gu{ખોપરી} (khoprī)}{ખોપરી (khoprī)}]{\gu{ખોપરી} (khoprī) — the skull}

\label{lesson:GU-C472-khopri}



\begin{quote}
Before the new one: say the Gujarati for a vein, then the Gujarati for a muscle.
\end{quote}

\subsection*{You'll want to know: \gu{ખોપરી}}
\textbf{\gu{ખોપરી}} — \emph{khoprī} — \textquotedblleft{}the skull\textquotedblright{}.

\begin{quote}
\gu{ખોપરી} \gu{મગજને} \gu{બચાવે} \gu{છે}.
\end{quote}

\begin{quote}
\emph{khoprī magajne bacāve che.}
\end{quote}

\begin{quote}
The skull protects the brain.
\end{quote}

A feminine noun.

\begin{cousinweb}
\textbf{\gu{ખોપરી}} (\emph{khoprī}) goes back to a reconstructed \emph{*khōppa-}, \textquotedblleft{}coconut shell\textquotedblright{}; CDIAL lists the Gujarati word, \textquotedblleft{}skull\textquotedblright{} (CDIAL 3936): the shell of the head.
\end{cousinweb}

\begin{grammarlens}[title={What you've built}]
One more word for describing what you lived and felt.
\end{grammarlens}

\subsection*{Your turn}
\begin{itemize}
\raggedright
  \item \emph{Say it:} \emph{khoprī}
  \item \emph{Say it:} \emph{khoprī}, once more
  \item \emph{Say it:} \emph{khoprī magajne bacāve che.} — The skull protects the brain.
  \item \emph{Recall:} say the Gujarati for a vein, then the Gujarati for a muscle, then say \emph{khoprī} again
\end{itemize}

\begin{tcolorbox}[breakable,colback=teal!4,colframe=teal!35!black,title={Before you move on}]
What does \gu{ખોપરી} mean? (\textquotedblleft{}The skull\textquotedblright{}.) Say it once more.
\end{tcolorbox}
\section[\texorpdfstring{\gu{આંતરડું} (ā̃tarḍũ)}{આંતરડું (ā̃tarḍũ)}]{\gu{આંતરડું} (ā̃tarḍũ) — an intestine}

\label{lesson:GU-C472-atardu}



\begin{quote}
Before the new one: say the Gujarati for a muscle, then the Gujarati for the skull.
\end{quote}

\subsection*{You'll want to know: \gu{આંતરડું}}
\textbf{\gu{આંતરડું}} — \emph{ā̃tarḍũ} — \textquotedblleft{}an intestine\textquotedblright{}.

\begin{quote}
\gu{વાસી} \gu{ભોજનથી} \gu{આંતરડામાં} \gu{ચેપ} \gu{લાગે} \gu{છે}.
\end{quote}

\begin{quote}
\emph{vāsī bhojanthī ā̃tarḍāmā̃ cep lāge che.}
\end{quote}

\begin{quote}
Stale food causes infection in the intestines.
\end{quote}

A neuter noun.

\begin{cousinweb}
\textbf{\gu{આંતરડું}} (\emph{ā̃tarḍũ}) is inherited from Sanskrit \emph{āntra}, \textquotedblleft{}entrail\textquotedblright{}; CDIAL lists the Gujarati word, \textquotedblleft{}entrail\textquotedblright{} (CDIAL 1182), with the noun ending \emph{-ḍũ}.
\end{cousinweb}

\begin{grammarlens}[title={What you've built}]
One more word for describing what you lived and felt.
\end{grammarlens}

\subsection*{Your turn}
\begin{itemize}
\raggedright
  \item \emph{Say it:} \emph{ā̃tarḍũ}
  \item \emph{Say it:} \emph{ā̃tarḍũ}, once more
  \item \emph{Say it:} \emph{vāsī bhojanthī ā̃tarḍāmā̃ cep lāge che.} — Stale food causes infection in the intestines.
  \item \emph{Recall:} say the Gujarati for a muscle, then the Gujarati for the skull, then say \emph{ā̃tarḍũ} again
\end{itemize}

\begin{tcolorbox}[breakable,colback=teal!4,colframe=teal!35!black,title={Before you move on}]
What does \gu{આંતરડું} mean? (\textquotedblleft{}An intestine\textquotedblright{}.) Say it once more.
\end{tcolorbox}
\section[\texorpdfstring{\gu{હોજરી} (hojrī)}{હોજરી (hojrī)}]{\gu{હોજરી} (hojrī) — the stomach (the organ)}

\label{lesson:GU-C472-hojri}



\begin{quote}
Before the new one: say the Gujarati for the skull, then the Gujarati for an intestine.
\end{quote}

\subsection*{You'll want to know: \gu{હોજરી}}
\textbf{\gu{હોજરી}} — \emph{hojrī} — \textquotedblleft{}the stomach (the organ)\textquotedblright{}.

\begin{quote}
\gu{વધારે} \gu{મરચાંથી} \gu{હોજરીમાં} \gu{બળે} \gu{છે}.
\end{quote}

\begin{quote}
\emph{vadhāre marcā̃thī hojrīmā̃ baḷe che.}
\end{quote}

\begin{quote}
Too many chillies make your stomach burn.
\end{quote}

A feminine noun.

\begin{cousinweb}
\textbf{\gu{હોજરી}} (\emph{hojrī}) is inherited: CDIAL lists the Gujarati word, \textquotedblleft{}stomach\textquotedblright{}, under Sanskrit \emph{ūbadhya}, \textquotedblleft{}undigested grass in the stomach of an animal\textquotedblright{} (CDIAL 2417).
\end{cousinweb}

\begin{grammarlens}[title={What you've built}]
One more word for describing what you lived and felt.
\end{grammarlens}

\subsection*{Your turn}
\begin{itemize}
\raggedright
  \item \emph{Say it:} \emph{hojrī}
  \item \emph{Say it:} \emph{hojrī}, once more
  \item \emph{Say it:} \emph{vadhāre marcā̃thī hojrīmā̃ baḷe che.} — Too many chillies make your stomach burn.
  \item \emph{Recall:} say the Gujarati for the skull, then the Gujarati for an intestine, then say \emph{hojrī} again
\end{itemize}

\begin{tcolorbox}[breakable,colback=teal!4,colframe=teal!35!black,title={Before you move on}]
What does \gu{હોજરી} mean? (\textquotedblleft{}The stomach (the organ)\textquotedblright{}.) Say it once more.
\end{tcolorbox}
